\documentclass[border=2pt, multi, tikz]{standalone}
\usepackage[T1]{fontenc}
\usepackage{lmodern}
\usepackage{amsmath}
\usetikzlibrary{arrows.meta,calc,positioning,backgrounds}
% Same-directory inputs so the file compiles both from the build script and from
% an editor that sets the working directory to this file's own folder.
% Shared colour palette for the Mean Teacher dissertation figures.
%
% This file used to carry drawing primitives as well (a perspective feature-volume
% macro, a compact backbone pic and four data-tile pics). Those are gone: the
% backbone panel now draws with the vendored PlotNeuralNet primitives, the protocol
% panel builds its network glyphs from that same primitive, and its data tiles are
% defined locally beside the drawing that uses them. What remains is the palette,
% which is the only thing genuinely shared.
%
% Role coding: teacher warm, student blue, labelled-data green, loss/topology
% purple. The probability ramp is a dark-to-yellow sequence so every predicted
% probability tile in the figure reads on the same scale.
\definecolor{MTInk}{HTML}{172033}
\definecolor{MTMuted}{HTML}{5F6B7A}
\definecolor{MTStudent}{HTML}{1464A5}
\definecolor{MTStudentLight}{HTML}{DCEEFE}
\definecolor{MTTeacher}{HTML}{D97706}
\definecolor{MTTeacherLight}{HTML}{FEEDD3}
\definecolor{MTLabel}{HTML}{11823B}
\definecolor{MTLabelLight}{HTML}{E1F5E8}
\definecolor{MTLoss}{HTML}{6D3FB5}
\definecolor{MTLossLight}{HTML}{EEE7FA}
\definecolor{MTProbabilityDark}{HTML}{260B45}
\definecolor{MTProbabilityMid}{HTML}{A33B78}
\definecolor{MTProbabilityHot}{HTML}{F59E0B}
\definecolor{MTProbabilityPeak}{HTML}{FFF29A}

% The network glyphs use the same vendored PlotNeuralNet primitive as the backbone
% panel, so the two figures read as one system at two zoom levels.
\input{plotneuralnet_layers.tex}

% Two-stream Mean Teacher protocol panel.
%
% Routing principle: the corridor between the teacher row and the student row is
% empty, so every operand of the consistency loss is delivered THROUGH it, and no
% wire travels further than one row.  Every segment is axis-parallel: each arrow
% either shares its source and target y (a single horizontal run) or turns through
% explicit right angles.  Rows are therefore placed so that a block's output
% anchor and its target's input anchor sit at exactly the same height.
\newcommand{\MTFlowBody}{\fontsize{12.7}{14.2}\selectfont}
\newcommand{\MTFlowSmall}{\fontsize{11.2}{12.5}\selectfont}
\newcommand{\MTFlowRole}{\fontsize{14.2}{15.8}\selectfont\bfseries}
\newcommand{\MTTeacherTarget}{T_U=\operatorname{sg}(p_{T,U})}

% Role fills for the network glyphs.  The teacher reuses the backbone's own
% convolution colours; the student is the matching blue, so the panel keeps its
% orange/blue role coding while sharing the backbone's shape grammar.
\def\MTTeacherFill{rgb:yellow,5;red,2.5;white,5}
\def\MTTeacherBand{rgb:yellow,5;red,5;white,5}
\def\MTStudentFill{rgb:blue,3;green,1;white,7}
\def\MTStudentBand{rgb:blue,4;green,2;white,4}

\tikzset{
  mt flow neutral/.style={-{Latex[length=1.85mm,width=1.22mm]},draw=MTMuted,line width=.70pt},
  mt flow topology/.style={-{Latex[length=1.85mm,width=1.22mm]},draw=MTLoss,line width=.74pt},
  mt flow supervised/.style={-{Latex[length=1.85mm,width=1.22mm]},draw=MTLabel,line width=.74pt},
  mt flow ema/.style={-{Latex[length=2.1mm,width=1.45mm]},draw=MTTeacher,line width=1.20pt,
    dash pattern=on 1.8mm off .85mm},
  mt flow grad/.style={-{Latex[length=2.0mm,width=1.35mm]},draw=MTStudent,line width=.86pt,
    dash pattern=on 1.7mm off .9mm},
  mt flow target/.style={draw=MTTeacher,fill=MTTeacherLight,rounded corners=.7mm,
    line width=.68pt,minimum width=2.62cm,minimum height=.70cm,inner sep=.7mm,
    font=\MTFlowBody,text=MTInk,align=center},
  mt flow operator/.style={circle,draw=MTLoss,fill=white,line width=.65pt,
    minimum size=.62cm,inner sep=0pt,font=\MTFlowBody,text=MTLoss},
  mt flow perturb/.style={circle,draw=MTStudent,fill=MTStudentLight,line width=.65pt,
    minimum size=.58cm,inner sep=0pt,font=\MTFlowBody,text=MTStudent},
  mt flow consistency/.style={draw=MTLoss,fill=MTLossLight,rounded corners=.8mm,
    line width=.78pt,text width=6.05cm,minimum height=2.95cm,inner sep=1.4mm,
    align=center,font=\MTFlowBody,text=MTInk},
  mt flow supervision/.style={draw=MTLabel,fill=MTLabelLight,rounded corners=.8mm,
    line width=.78pt,minimum width=3.70cm,minimum height=1.25cm,inner sep=1.0mm,
    align=center,font=\MTFlowBody,text=MTInk},
  mt flow total/.style={draw=MTInk,fill=white,rounded corners=.8mm,line width=.88pt,
    minimum width=3.00cm,minimum height=1.22cm,inner sep=1.0mm,align=center,
    font=\MTFlowBody,text=MTInk},
}

% Five-stage symbolic backbone built from the backbone panel's own RightBandedBox
% primitive: same cuboid shape, same banded convolution grammar, only the role
% colour and the scale differ.  It summarises rather than enumerates the exact
% 11-stage network; the companion figure carries the real stage list.
% #1 name, #2 x, #3 flow-axis y, #4 fill colour, #5 band colour.
\newcommand{\MTUNet}[5]{%
  \pic at (#2,#3) {RightBandedBox={name=#1s1,fill=#4,bandfill=#5,height=9,width=2,depth=6}};
  \pic[shift={(.18,0,0)}] at (#1s1-east) {RightBandedBox={name=#1s2,fill=#4,bandfill=#5,height=7,width=2.5,depth=5}};
  \pic[shift={(.18,0,0)}] at (#1s2-east) {RightBandedBox={name=#1s3,fill=#4,bandfill=#5,height=5,width=3,depth=4}};
  \pic[shift={(.18,0,0)}] at (#1s3-east) {RightBandedBox={name=#1s4,fill=#4,bandfill=#5,height=7,width=2.5,depth=5}};
  \pic[shift={(.18,0,0)}] at (#1s4-east) {RightBandedBox={name=#1s5,fill=#4,bandfill=#5,height=9,width=2,depth=6}};
  \foreach \i/\j in {1/2,2/3,3/4,4/5}{%
    \draw[-{Latex[length=1.05mm,width=.72mm]},draw=MTInk!55,line width=.42pt]
      (#1s\i-east) -- (#1s\j-west);}
  % Incoming ports are pulled 0.27 left of the pic's own -west anchor. That anchor
  % lies in the z=0 plane, but the drawn front face is projected down-left by
  % 0.385*depth/2, so an arrow ending on the raw anchor puts its head INSIDE the
  % visible block. The offset is x-only, so the runs stay horizontal.
  \coordinate (#1-west)  at ($(#1s1-west)+(-.27,0)$);
  \coordinate (#1-westu) at ($(#1s1-west)+(-.27,.25)$);
  \coordinate (#1-westl) at ($(#1s1-west)+(-.27,-.25)$);
  \coordinate (#1-east)  at (#1s5-east);
  \coordinate (#1-eastu) at ($(#1s5-east)+(0,.25)$);
  \coordinate (#1-eastl) at ($(#1s5-east)+(0,-.25)$);
  % The waist of the U is the only clear vertical corridor through the glyph.
  % Same projection correction, vertically: these are arrowhead targets for the
  % EMA and gradient routes, so they must clear the drawn faces.
  \coordinate (#1-waistn) at ($(#1s3-north)+(0,.20)$);
  \coordinate (#1-waists) at ($(#1s3-south)+(0,-.20)$);
}

% Local input / output tiles retain the established semantics.
\tikzset{pics/mt flow ct/.style args={#1}{code={
  \path[draw=MTMuted!65,fill=black!7,line width=.55pt] (0,0) rectangle (1.05,.86);
  \path[draw=MTMuted!65,fill=black!2,line width=.55pt]
    (0,.86) -- (.16,1.00) -- (1.21,1.00) -- (1.05,.86) -- cycle;
  \path[draw=MTMuted!65,fill=black!18,line width=.55pt]
    (1.05,0) -- (1.21,.14) -- (1.21,1.00) -- (1.05,.86) -- cycle;
  \fill[black!72] (.36,.46) ellipse (.20 and .29);
  \fill[black!72] (.70,.46) ellipse (.20 and .29);
  \draw[white!84,line width=.52pt,line cap=round] (.525,.69) -- (.525,.51) -- (.42,.39) (.525,.51) -- (.64,.39);
  \node[font=\MTFlowBody,text=MTInk] at (.525,-.27) {#1};
  \coordinate (-west) at (0,.50); \coordinate (-east) at (1.21,.50);
  \coordinate (-north) at (.605,1.00); \coordinate (-south) at (.525,0);
}}}

\tikzset{pics/mt flow probability/.style args={#1}{code={
  \path[draw=MTProbabilityMid!55,fill=MTProbabilityMid!14,line width=.45pt] (.08,.08) rectangle (1.06,.88);
  \path[draw=MTProbabilityDark!90,fill=MTProbabilityDark,line width=.58pt] (0,0) rectangle (.98,.80);
  \draw[MTProbabilityMid,line width=4.1pt,line cap=round,line join=round,opacity=.66]
    (.49,.69) -- (.49,.50) -- (.32,.35) -- (.22,.17) (.49,.50) -- (.67,.35) -- (.81,.17);
  \draw[MTProbabilityHot,line width=2.25pt,line cap=round,line join=round]
    (.49,.69) -- (.49,.50) -- (.32,.35) -- (.22,.17) (.32,.35) -- (.16,.30)
    (.49,.50) -- (.67,.35) -- (.81,.17) (.67,.35) -- (.86,.30);
  \draw[MTProbabilityPeak,line width=.72pt,line cap=round,line join=round]
    (.49,.69) -- (.49,.50) -- (.32,.35) -- (.22,.17) (.49,.50) -- (.67,.35) -- (.81,.17);
  \node[font=\MTFlowBody,text=MTInk] at (.49,-.27) {#1};
  \coordinate (-west) at (0,.40); \coordinate (-east) at (1.06,.40);
  \coordinate (-north) at (.53,.88); \coordinate (-south) at (.49,0);
}}}

\tikzset{pics/mt flow skeleton/.style args={#1}{code={
  \path[draw=MTLoss!28,fill=MTLoss!4,line width=.45pt] (.08,.08) rectangle (1.02,.88);
  \path[draw=MTLoss!72,fill=white,line width=.58pt] (0,0) rectangle (.94,.80);
  \draw[MTLoss!26,line width=2.9pt,line cap=round,line join=round]
    (.47,.69) -- (.47,.50) -- (.31,.35) -- (.21,.17) (.47,.50) -- (.64,.35) -- (.78,.17);
  \draw[MTLoss,line width=1.02pt,line cap=round,line join=round]
    (.47,.69) -- (.47,.50) -- (.31,.35) -- (.21,.17) (.31,.35) -- (.15,.30)
    (.47,.50) -- (.64,.35) -- (.78,.17) (.64,.35) -- (.84,.30);
  \node[font=\MTFlowBody,text=MTInk] at (.47,-.27) {#1};
  \coordinate (-west) at (0,.40); \coordinate (-east) at (1.02,.40);
  \coordinate (-north) at (.51,.88); \coordinate (-south) at (.47,0);
}}}

\tikzset{pics/mt flow mask/.style args={#1}{code={
  \path[draw=MTLabel!28,fill=MTLabel!5,line width=.45pt] (.08,.08) rectangle (1.02,.88);
  \path[draw=MTLabel!72,fill=MTLabelLight,line width=.58pt] (0,0) rectangle (.94,.80);
  \draw[MTLabel,line width=1.75pt,line cap=round,line join=round]
    (.47,.69) -- (.47,.50) -- (.31,.35) -- (.21,.17) (.31,.35) -- (.15,.30)
    (.47,.50) -- (.64,.35) -- (.78,.17) (.64,.35) -- (.84,.30);
  \node[font=\MTFlowBody,text=MTInk] at (.47,-.27) {#1};
  \coordinate (-west) at (0,.40); \coordinate (-east) at (1.02,.40);
  \coordinate (-north) at (.51,.88); \coordinate (-south) at (.47,0);
}}}

\begin{document}
\begin{tikzpicture}[x=1cm,y=1cm]
  \path[use as bounding box] (.00,-.55) rectangle (28.90,9.55);

  % ---------------------------------------------------------------------------
  % Row heights.  Teacher 7.70, student 4.10; the student's two ports sit at
  % 4.35 (unlabelled) and 3.85 (labelled).  Every downstream tile is placed so
  % its anchor lands on exactly one of these lines.
  % ---------------------------------------------------------------------------

  % Role banners stay short: the dataset marking rides on the input tiles instead,
  % which is where the sampling actually happens and keeps this column clear of the
  % perturbation blocks on the risers.
  \node[anchor=south west,font=\MTFlowRole,text=MTInk] at (.20,8.32) {Unlabelled};
  \node[anchor=south west,font=\MTFlowRole,text=MTInk] at (.20,2.32) {Labelled};

  \pic (xu) at (.28,7.20) {mt flow ct={$x_U\in\mathcal D_U$}};
  \pic (xl) at (.28,1.20) {mt flow ct={$x_L\in\mathcal D_L$}};

  \MTUNet{tch}{4.10}{7.70}{\MTTeacherFill}{\MTTeacherBand}
  \MTUNet{stu}{4.10}{4.10}{\MTStudentFill}{\MTStudentBand}
  % Role labels start left of their glyphs so the student's cannot be crossed by
  % the EMA riser running up the waist at x = 5.66.
  \node[anchor=south west,font=\MTFlowBody\bfseries,text=MTTeacher] at (2.90,8.95)
    {Teacher $f_{\theta_T}$};
  \node[anchor=south west,font=\MTFlowBody\bfseries,text=MTStudent] at (2.90,5.35)
    {Student $f_{\theta_S}$};

  % Data routes.  One x_U sample forks to both networks; the labelled riser is
  % held at x=3.00 and the unlabelled fork at x=2.30 so the two never cross.
  \coordinate (ufork) at (2.30,7.70);
  \fill[MTMuted] (ufork) circle (1.15pt);
  \draw[mt flow neutral] (xu-east) -- (tch-west);
  \draw[mt flow neutral] (ufork) -- (2.30,4.35) -- (stu-westu);
  \draw[mt flow neutral] (xl-east) -- (3.00,1.70) -- (3.00,3.85) -- (stu-westl);

  % The student's additional intensity perturbation, applied to its whole input.
  % Labelled with the operator itself rather than with eta alone: the transform is a
  % per-sample scale and shift plus per-voxel noise, not noise only.
  % Collapsed to the figure's existing operator grammar: a circle is an operator,
  % so the perturbation is one glyph and the transform itself lives in the caption.
  % pi, not a tilde: the tilde already means "distributed as" in the same equation.
  \node[mt flow perturb] (pertu) at (2.30,6.05) {$\pi$};
  \node[mt flow perturb] (pertl) at (3.00,2.65) {$\pi$};
  \node[anchor=west,font=\MTFlowSmall,text=MTStudent] at (2.42,5.30) {$\tilde x_U$};
  \node[anchor=west,font=\MTFlowSmall,text=MTStudent] at (3.12,3.30) {$\tilde x_L$};

  % EMA rises through the waist of the two U silhouettes, the one clear corridor.
  \draw[mt flow ema] (stu-waistn) --
    node[pos=.42,right=1.2mm,font=\MTFlowSmall,text=MTTeacher,fill=white,inner sep=.35mm] {EMA}
    (tch-waists);

  % Prediction column.  Each tile's west anchor is on its source's output line, so
  % these are single horizontal runs.
  \pic (pt)  at (9.60,7.30) {mt flow probability={$p_{T,U}$}};
  \pic (psu) at (9.60,3.95) {mt flow probability={$p_{S,U}$}};
  \pic (psl) at (9.60,1.60) {mt flow probability={$p_{S,L}$}};
  \draw[mt flow neutral] (tch-east)  -- (pt-west);
  \draw[mt flow neutral] (stu-eastu) -- (psu-west);
  \draw[mt flow neutral] (stu-eastl) -- (8.60,3.85) -- (8.60,2.00) -- (psl-west);

  % Teacher target and the two soft-skeleton operators, all on the row lines.
  % Routing the consistency loss into the TOP of the total-loss node frees the
  % corridor that used to run between them, and that width is spent here: the
  % target, operator and skeleton tile each get roughly 0.75 cm of clear air.
  \node[mt flow target]   (target) at (12.55,7.70) {$\MTTeacherTarget$};
  \node[mt flow operator] (skt)    at (15.00,7.70) {$S$};
  \node[mt flow operator] (sks)    at (15.00,4.35) {$S$};
  \pic (kt) at (16.10,7.30) {mt flow skeleton={$S(T_U)$}};
  \pic (ks) at (16.10,3.95) {mt flow skeleton={$S(p_{S,U})$}};
  \draw[mt flow neutral]  (pt-east) -- (target.west);
  \draw[mt flow topology] (target.east) -- (skt.west);
  \draw[mt flow topology] (skt.east) -- (kt-west);
  \draw[mt flow topology] (psu-east) -- (sks.west);
  \draw[mt flow topology] (sks.east) -- (ks-west);

  % Consistency loss: four operands entering on four ports spaced symmetrically
  % about the box's centre line, teacher-derived above, student-derived below.
  \node[mt flow consistency] (lcons) at (21.30,6.00)
    {\textbf{soft clDice consistency}\\[1.1mm]
     $T_{\mathrm{prec}}=\dfrac{\sum S(p_{S,U})\,T_U+s}{\sum S(p_{S,U})+s}$\\[1.6mm]
     $T_{\mathrm{sens}}=\dfrac{\sum S(T_U)\,p_{S,U}+s}{\sum S(T_U)+s}$\\[1.9mm]
     $\mathcal L_{\mathrm{cons}}
       =1-\dfrac{2\,T_{\mathrm{prec}}T_{\mathrm{sens}}}
                {T_{\mathrm{prec}}+T_{\mathrm{sens}}+\varepsilon}$};
  \coordinate (lcons-kt)  at ([yshift= 8.5mm]lcons.west);
  \coordinate (lcons-tu)  at ([yshift= 3.5mm]lcons.west);
  \coordinate (lcons-psu) at ([yshift=-3.5mm]lcons.west);
  \coordinate (lcons-ks)  at ([yshift=-8.5mm]lcons.west);

  \draw[mt flow topology] (kt-east) -- (17.75,7.70) -- (17.75,6.85) -- (lcons-kt);
  \draw[mt flow topology] (target.south) -- (12.55,6.35) -- (lcons-tu);
  \draw[mt flow topology] (11.40,4.35) -- (11.40,5.65) -- (lcons-psu);
  \fill[MTLoss] (11.40,4.35) circle (1.15pt);
  \draw[mt flow topology] (ks-east) -- (17.75,4.35) -- (17.75,5.15) -- (lcons-ks);

  % Labelled supervision.  y_L sits clear of the p_{S,L} run rather than on it,
  % and both operands enter on ports spaced symmetrically about the box centre.
  \pic (yl) at (11.60,0.70) {mt flow mask={$y_L$}};
  \node[mt flow supervision] (lsup) at (18.85,1.55)
    {\textbf{Dice + cross-entropy}\\[-.1mm]
     $\mathcal L_{\mathrm{sup}}(p_{S,L},y_L)$};
  \coordinate (lsup-p) at ([yshift= 4.5mm]lsup.west);
  \coordinate (lsup-y) at ([yshift=-4.5mm]lsup.west);
  \draw[mt flow supervised] (psl-east) -- (lsup-p);
  \draw[mt flow supervised] (yl-east)  -- (lsup-y);

  % Total objective.  Moved right to clear the widened consistency block.
  \node[mt flow total] (total) at (27.00,4.00)
    {$\mathcal L=\mathcal L_{\mathrm{sup}}$\\
     $+\,w(e)\,\mathcal L_{\mathrm{cons}}$};
  % The two losses enter on different faces: consistency from above, supervision
  % from the west, so they need no shared vertical corridor.
  \draw[mt flow topology]   (lcons.east) -- (27.00,6.00) -- (total.north);
  \draw[mt flow supervised] (lsup.east)  -- (24.90,1.55) -- (24.90,4.00) -- (total.west);

  % The gradient returns to the student and to nothing else, so it is drawn as a
  % closed route rather than a dangling stub with a caption on it.
  \draw[mt flow grad] (total.south) -- (27.00,-.15) -- (5.66,-.15) -- (stu-waists);
  \node[font=\MTFlowSmall,text=MTStudent,fill=white,inner xsep=1.0mm,inner ysep=.4mm]
    at (15.20,-.15) {gradient to $\theta_S$ only};
\end{tikzpicture}
\end{document}
